\Problem{Cheapest Routes}
\Statement
We have collected abundant
information about the local roads and accommodations in a region that
we will traverse. Our plan is to go from city $A$ to city $B$ and
we would like to spend the least possible money. For each road
connecting two cities $u$ and $v$ we know the cost $\omega(u,v)=\omega(v,u)$
to travel
along that road (tolls, fuel, meals during the journey, \ldots). Every
time we go from a city $u$ to one of its neighbors $v$ we must stop at
$v$ and spend there one night; we know the cost $\omega'(v)$
of stopping at each city $v$ (the cost added by $A$ and $B$ to
our route is 0, since they
are our initial and final points).
All costs, of vertices and of edges, are positive. Thus the cost
of the route 
\[
P = [A, v_1, \ldots, v_n, B]
\]
is
\[
  \text{cost}(P) = \omega(A,v_1) + \omega(v_1,v_2) + \ldots + \omega(v_n,B) 
  + \omega'(v_1)+\ldots+\omega'(v_n).
\]

Write a program in C++ which, given an undirected weighted graph with
positive
costs at the vertices and at the edges, and two vertices $A$ and $B$,
returns the cost of the cheapest route to go from $A$ to $B$, or an
indication that not such route exists.
%% If several different paths
%% between $A$ and $B$ have minimal cost, return the one with less edges;
%% if there are ties, return the one for which the sequence of vertices
%% is lexicographically smallest.

\Input
All data in the input are positive integers.
The input starts with two integers $2{\le}n{\le}10000$ and $m$, $0{\le}m\le{20n}$.
After that, a sequence of positive
integers $\omega'(0), \ldots, \omega'(n-1)$ of
the weights $\omega'(u)$ of the
$n$ vertices of the graph.  Then the input contains
a sequence of the $m$ edges in the graph
as triplets of the form ${\langle}u,v,\omega(u,v)\rangle$. Vertices
$u$ and $v$ are integers 
in $\{0,\ldots,n-1\}$ and the weights $\omega(u,v)$ are positive integers.
No weight, of edges or vertices, is larger than 100000.

You can assume that there are no two different edges connecting
the same pair of vertices nor any edge connecting a vertex to itself.
Finally, there is
a sequence of pairs ${\langle}A_i, B_i\rangle$, with each $A_i$ and $B_i$
denoting vertices of the graph ($0{\le}A_i,B_i<n$).

\Output
For each pair ${\langle}A_i, B_i\rangle$ in the input sequence the program
writes the cost $\delta$ of the cheapest route between $A_i$ and $B_i$. with the format
\texttt{c($A_i$,$B_i$) = $\delta$}.
% and the cheapest path itself
% (a sequence of vertices starting
% with $A_i$ and ending with $B_i$)
If no route exists between $A_i$ and $B_i$ the program
writes \texttt{c($A_i$,$B_i$) = +oo}. The ouput for each case is
ended with a newline (\texttt{endl}).

\SampleTwoCol
